% Options for packages loaded elsewhere
\PassOptionsToPackage{unicode}{hyperref}
\PassOptionsToPackage{hyphens}{url}
\documentclass[
]{article}
\usepackage{xcolor}
\usepackage{amsmath,amssymb}
\setcounter{secnumdepth}{-\maxdimen} % remove section numbering
\usepackage{iftex}
\ifPDFTeX
  \usepackage[T1]{fontenc}
  \usepackage[utf8]{inputenc}
  \usepackage{textcomp} % provide euro and other symbols
\else % if luatex or xetex
  \usepackage{unicode-math} % this also loads fontspec
  \defaultfontfeatures{Scale=MatchLowercase}
  \defaultfontfeatures[\rmfamily]{Ligatures=TeX,Scale=1}
\fi
\usepackage{lmodern}
\ifPDFTeX\else
  % xetex/luatex font selection
\fi
% Use upquote if available, for straight quotes in verbatim environments
\IfFileExists{upquote.sty}{\usepackage{upquote}}{}
\IfFileExists{microtype.sty}{% use microtype if available
  \usepackage[]{microtype}
  \UseMicrotypeSet[protrusion]{basicmath} % disable protrusion for tt fonts
}{}
\makeatletter
\@ifundefined{KOMAClassName}{% if non-KOMA class
  \IfFileExists{parskip.sty}{%
    \usepackage{parskip}
  }{% else
    \setlength{\parindent}{0pt}
    \setlength{\parskip}{6pt plus 2pt minus 1pt}}
}{% if KOMA class
  \KOMAoptions{parskip=half}}
\makeatother
\usepackage{longtable,booktabs,array}
\usepackage{caption}
\captionsetup[table]{skip=6pt}
\newcounter{none} % for unnumbered tables
\usepackage{calc} % for calculating minipage widths
% Correct order of tables after \paragraph or \subparagraph
\usepackage{etoolbox}
\makeatletter
\patchcmd\longtable{\par}{\if@noskipsec\mbox{}\fi\par}{}{}
\makeatother
\setlength{\emergencystretch}{3em} % prevent overfull lines
\providecommand{\tightlist}{%
  \setlength{\itemsep}{0pt}\setlength{\parskip}{0pt}}
\usepackage{bookmark}
\IfFileExists{xurl.sty}{\usepackage{xurl}}{} % add URL line breaks if available
\urlstyle{same}
% fallback for those not using the hyperref driver hyperxmp:
\makeatletter
\@ifundefined{xmpquote}{\newcommand{\xmpquote}[1]{#1}}{}
\makeatother
\hypersetup{
  hidelinks,
  pdfcreator={LaTeX via pandoc}}

\author{}
\date{}

\begin{document}

\section{Cubli: study guide and first math
problem}\label{cubli-study-guide-and-first-math-problem}

\subsection{1. What the ETH reference
says}\label{what-the-eth-reference-says}

\textbf{Michael Muehlebach (2018), \emph{On the Stabilization of
Unstable Systems - Some Comments and Applications}, ETH thesis no.
24789.} \href{phdthesisE-Collection.pdf}{Local PDF} \textbar{}
\href{https://doi.org/10.3929/ethz-b-000254022}{Official record}.

The thesis has four parts: Cubli, a ducted-fan flying vehicle,
constrained LQR/MPC approximations, and distributed state estimation.
\textbf{Read Part A first}; the other parts are not prerequisites for
building a balancing cube. Page numbers below are printed thesis pages;
PDF viewer pages are also given.

{\def\LTcaptype{none} % do not increment counter
\begin{longtable}[]{@{}
  >{\raggedright\arraybackslash}p{(\linewidth - 4\tabcolsep) * \real{0.3333}}
  >{\raggedright\arraybackslash}p{(\linewidth - 4\tabcolsep) * \real{0.3333}}
  >{\raggedright\arraybackslash}p{(\linewidth - 4\tabcolsep) * \real{0.3333}}@{}}
\toprule\noalign{}
\begin{minipage}[b]{\linewidth}\raggedright
Read
\end{minipage} & \begin{minipage}[b]{\linewidth}\raggedright
Main result
\end{minipage} & \begin{minipage}[b]{\linewidth}\raggedright
What it means for this project
\end{minipage} \\
\midrule\noalign{}
\endhead
\bottomrule\noalign{}
\endlastfoot
P1, dynamics: pp.~25-28 (PDF 44-47) & A 150 mm cube with three
perpendicular reaction wheels; gravity, inertia, and angular momentum
determine its motion. & Motor torque accelerates a wheel and reacts on
the frame. Wheel speeds are states, not just outputs. \\
P1, nonlinear control: pp.~28-37 (PDF 47-56) & Backstepping with a
Lyapunov argument stabilizes the upright set almost everywhere in the
ideal model. Feedback linearization supports trajectory tracking. &
First understand stability proofs and the model; the result is not a
guarantee under arbitrary motor limits or impacts. \\
P1, jump-up: pp.~37-45 (PDF 56-64) & Braking transfers wheel momentum to
the cube. A gradient-based learning method adjusts initial wheel speeds;
a guiding controller tracks the rise. & Jump-up is a separate hybrid
maneuver, not ordinary balancing with larger gains. \\
P1, experiments: pp.~45-47 (PDF 64-66) & Reported steady-state RMS
inclination error below 0.025 degrees; recovery in under 1.8 s after a
0.17 N m, 60 ms disturbance on one wheel. & These are results for ETH's
apparatus, not performance promises for a printed prototype. \\
P2, estimation: pp.~51-78 (PDF 70-97) & Multiple accelerometers and a
non-accelerated pivot assumption estimate tilt using rigid-body
kinematics, without mass/inertia knowledge. & A moving accelerometer
measures rotational acceleration as well as gravity; simply treating its
reading as gravity can give wrong tilt. \\
\end{longtable}
}

\textbf{Important limitations:} the reduced gravity-vector
representation describes tilt but discards yaw orientation. In the ideal
frictionless-pivot model, angular momentum about gravity is conserved:
one cannot generally force both frame yaw rate and all wheel speeds to
zero. Ground friction changes this. ETH's almost-global balancing result
concerns an upright set, not arbitrary full-attitude regulation with
every wheel stopped.

\textbf{INDI is not the controller developed in these Cubli chapters.}
It is our proposed comparison. Start with their dynamics and a
conventional baseline; do not assume INDI will outperform it without
experiments.

\subsection{2. Nonlinear-control book}\label{nonlinear-control-book}

\textbf{Hassan K. Khalil, \emph{Nonlinear Systems}, 3rd edition,
Prentice Hall, 2002, ISBN 9780130673893.}
\href{https://www.egr.msu.edu/~khalil/NonlinearSystems/index.html}{Author's
companion site} \textbar{}
\href{https://www.egr.msu.edu/~khalil/NonlinearSystems/contents.html}{Verified
contents} \textbar{}
\href{https://www.pearson.com/en-us/subject-catalog/p/nonlinear-systems/P200000003306/9780130673893}{Publisher}.
Obtain the book through a library or publisher; only its citation and
links are included here, not a copied textbook.

Suggested order: section 1.2.1 (pendulum), chapter 2 (phase portraits),
sections 4.1-4.3 (Lyapunov, invariance, linearization), section 12.2
(stabilization via linearization), chapter 13 (feedback linearization),
then section 14.3 (backstepping). Chapters 4 and 13 are the priority.

\subsection{3. Similar concept: a reaction-wheel inverted
pendulum}\label{similar-concept-a-reaction-wheel-inverted-pendulum}

Think of one Cubli face on a fixed horizontal pivot, carrying one motor
and wheel. This has the same gravity-versus-reaction-torque mechanism,
without three-dimensional coupling. The following is an original
educational exercise, not a reproduction of a textbook problem or ETH's
full model.

Assume a rigid body, fixed frictionless pivot, ideal motor torque, no
impacts, and no wheel/body drag. Let:

\begin{itemize}
\tightlist
\item
  \(\theta\): body angle from upright; \(\omega=\dot\theta\).
\item
  \(\phi\): wheel angle relative to the body; \(s=\dot\theta+\dot\phi\)
  is its absolute speed.
\item
  \(I\): pivot inertia of the non-spinning assembly, including
  wheel-center translation but excluding the wheel's axial spin inertia
  \(J\).
\item
  \(m\): total mass; \(l\): pivot-to-total-center-of-mass distance.
\item
  \(\tau\): torque on the wheel; the reaction torque on the body is
  \(-\tau\).
\end{itemize}

Use these \textbf{illustrative parameters, not a hardware shopping
specification}:

\[m=0.40\ {\rm kg},\quad l=0.08\ {\rm m},\quad
I=0.004\ {\rm kg\,m^2},\quad J=0.0002\ {\rm kg\,m^2},\quad g=9.81\ {\rm m/s^2}.\]

\subsubsection{Problems to solve}\label{problems-to-solve}

\begin{enumerate}
\def\labelenumi{\arabic{enumi}.}
\tightlist
\item
  Start with
  \(L=\frac12 I\dot\theta^2+\frac12 J(\dot\theta+\dot\phi)^2-mgl\cos\theta\).
  The generalized forces for \((\theta,\phi)\) are \((0,\tau)\) because
  the motor is internal. Derive the two Euler-Lagrange equations and
  show \(I\dot\omega=mgl\sin\theta-\tau\) and \(J\dot s=\tau\).
\item
  Linearize at \((\theta,\omega,s,\tau)=(0,0,0,0)\). Write \(A,B\) for
  \(x=(\theta,\omega,s)\); show the upright angle is open-loop unstable.
\item
  For the \textbf{angle subsystem}, choose
  \(\tau=mgl\sin\theta+I(k_p\theta+k_d\omega)\). Find the closed-loop
  equation and choose gains for a critically damped response with
  natural frequency \(8\ {\rm rad/s}\).
\item
  Use \(V=\frac12 k_p\theta^2+\frac12\omega^2\) to prove angle
  convergence. Explain why \(\dot V\leq0\) alone is insufficient and use
  LaSalle's principle. Does this also prove \(s\to0\)?
\item
  Starting from \(\theta=5\) degrees, \(\omega=s=0\), calculate the
  initial torque. Simulate the nonlinear plant with
  \(|\tau|\leq0.20\ {\rm N\,m}\) and a wheel-speed ceiling
  \(|s|\leq300\ {\rm rad/s}\). Detect and stop/report a speed-limit
  violation; do not clip \(s\), since that would introduce an unmodeled
  braking torque.
\item
  Write \(\dot\omega=f(\theta)+G\tau\) and derive an INDI update for
  target acceleration \(v=-k_p\theta-k_d\omega\), using a recent
  measured acceleration \(\alpha_0\) and the corresponding applied
  torque \(\tau_0\). Explain what fails if the acceleration measurement
  is delayed or noisy.
\end{enumerate}

Checkpoints: open after attempting the problems

\begin{itemize}
\tightlist
\item
  For the Euler-Lagrange coordinates, \(q_1=\theta\) (body angle) and
  \(q_2=\phi\) (wheel angle relative to the body). Their generalized
  forces are \(Q_1=Q_\theta=0\) and \(Q_2=Q_\phi=\tau\): the motor
  applies \(-\tau\) to the body and \(+\tau\) to the wheel, so its
  virtual work is
  \(-\tau\,\delta\theta+\tau\,\delta(\theta+\phi)=\tau\,\delta\phi\).
  Since
  \(\partial L/\partial\dot\theta=I\dot\theta+J(\dot\theta+\dot\phi)\)
  and \(\partial L/\partial\theta=mgl\sin\theta\), the \(\theta\)
  equation is \((I+J)\ddot\theta+J\ddot\phi-mgl\sin\theta=0\), or
  \((I+J)\ddot\theta+J\ddot\phi=mgl\sin\theta\). The \(\phi\) equation
  is \(J(\ddot\theta+\ddot\phi)=\tau\), i.e.~\(J\dot s=\tau\).
  Substituting this into the \(\theta\) equation gives
  \(I\ddot\theta=mgl\sin\theta-\tau\), or
  \(I\dot\omega=mgl\sin\theta-\tau\).
\item
  The coupled equations are
  \((I+J)\ddot\theta+J\ddot\phi=mgl\sin\theta\) and
  \(J(\ddot\theta+\ddot\phi)=\tau\).
\item
  With \(a=mgl/I=78.48\ {\rm s^{-2}}\):
  \(A=\begin{bmatrix}0&1&0\\a&0&0\\0&0&0\end{bmatrix}\),
  \(B=\begin{bmatrix}0\\-1/I\\1/J\end{bmatrix}\). The open-loop
  eigenvalues are \(0\) and approximately \(\pm8.86\ {\rm s^{-1}}\).
\item
  Ideal cancellation gives \(\ddot\theta+k_d\dot\theta+k_p\theta=0\);
  choose \(k_p=64\ {\rm s^{-2}}\), \(k_d=16\ {\rm s^{-1}}\).
\item
  \(\dot V=-k_d\omega^2\). The largest invariant set with \(\omega=0\)
  also has \(\theta=0\), proving angle convergence in the ideal,
  unsaturated angle model. This is not a global physical Cubli proof.
\item
  Wheel speed need not converge to zero:
  \(s(t)=s(0)+J^{-1}\int_0^t\tau\,dt\). The angle-only controller leaves
  the wheel's final speed unconstrained. As an extension, use full-state
  LQR to include wheel-speed regulation.
\item
  The initial torque is approximately \(0.0497\ {\rm N\,m}\). Even when
  initial torque is feasible, check the entire trajectory's limits.
\item
  \(f=(mgl/I)\sin\theta\), \(G=-1/I\). Locally,
  \(\dot\omega\approx\alpha_0+G(\tau-\tau_0)\), hence
  \(\tau=\tau_0+I(\alpha_0-v)\). This assumes gravity changes little
  over the increment and applied torque responds fast enough.
  Time-align/filter the torque and acceleration signals consistently and
  include actuator dynamics and saturation.
\end{itemize}

\end{document}
